\RequirePackage[abspath]{currfile}

%%
%% This is file `sample-manuscript.tex',
%% generated with the docstrip utility.
%%
%% The original source files were:
%%
%% samples.dtx  (with options: `all,proceedings,bibtex,manuscript')
%% 
%% IMPORTANT NOTICE:
%% 
%% For the copyright see the source file.
%% 
%% Any modified versions of this file must be renamed
%% with new filenames distinct from sample-manuscript.tex.
%% 
%% For distribution of the original source see the terms
%% for copying and modification in the file samples.dtx.
%% 
%% This generated file may be distributed as long as the
%% original source files, as listed above, are part of the
%% same distribution. (The sources need not necessarily be
%% in the same archive or directory.)
%%
%%
%% Commands for TeXCount
%TC:macro \cite [option:text,text]
%TC:macro \citep [option:text,text]
%TC:macro \citet [option:text,text]
%TC:envir table 0 1
%TC:envir table* 0 1
%TC:envir tabular [ignore] word
%TC:envir displaymath 0 word
%TC:envir math 0 word
%TC:envir comment 0 0
%%
%% The first command in your LaTeX source must be the \documentclass
%% command.
%%
%% For submission and review of your manuscript please change the
%% command to \documentclass[manuscript, screen, review]{acmart}.
%%
%% When submitting camera ready or to TAPS, please change the command
%% to \documentclass[sigconf]{acmart} or whichever template is required
%% for your publication.
%%
%%
\documentclass[manuscript,anonymous,review,authorversion=true,nonacm=true,screen]{acmart}
%%
%% \BibTeX command to typeset BibTeX logo in the docs
\AtBeginDocument{%
  \providecommand\BibTeX{{%
    Bib\TeX}}}

\usepackage{verbatimbox}

\usepackage{etoolbox}
\newtoggle{showoutline}
\togglefalse{showoutline}
\toggletrue{showoutline}

\newcommand{\outline}[1]{\iftoggle{showoutline}{#1}{}}

%% Rights management information.  This information is sent to you
%% when you complete the rights form.  These commands have SAMPLE
%% values in them; it is your responsibility as an author to replace
%% the commands and values with those provided to you when you
%% complete the rights form.
% \setcopyright{acmlicensed}
% \copyrightyear{2018}
% \acmYear{2018}
% \acmDOI{XXXXXXX.XXXXXXX}
%% These commands are for a PROCEEDINGS abstract or paper.
% \acmConference[Conference acronym 'XX]{Make sure to enter the correct
%   conference title from your rights confirmation email}{June 03--05,
%   2018}{Woodstock, NY}
%%
%%  Uncomment \acmBooktitle if the title of the proceedings is different
%%  from ``Proceedings of ...''!
%%
%%\acmBooktitle{Woodstock '18: ACM Symposium on Neural Gaze Detection,
%%  June 03--05, 2018, Woodstock, NY}
% \acmISBN{978-1-4503-XXXX-X/2018/06}


%%
%% Submission ID.
%% Use this when submitting an article to a sponsored event. You'll
%% receive a unique submission ID from the organizers
%% of the event, and this ID should be used as the parameter to this command.
%%\acmSubmissionID{123-A56-BU3}

%%
%% For managing citations, it is recommended to use bibliography
%% files in BibTeX format.
%%
%% You can then either use BibTeX with the ACM-Reference-Format style,
%% or BibLaTeX with the acmnumeric or acmauthoryear sytles, that include
%% support for advanced citation of software artefact from the
%% biblatex-software package, also separately available on CTAN.
%%
%% Look at the sample-*-biblatex.tex files for templates showcasing
%% the biblatex styles.
%%

%%
%% The majority of ACM publications use numbered citations and
%% references.  The command \citestyle{authoryear} switches to the
%% "author year" style.
%%
%% If you are preparing content for an event
%% sponsored by ACM SIGGRAPH, you must use the "author year" style of
%% citations and references.
%% Uncommenting
%% the next command will enable that style.
%%\citestyle{acmauthoryear}


%%
%% end of the preamble, start of the body of the document source.
\begin{document}

%%
%% The "title" command has an optional parameter,
%% allowing the author to define a "short title" to be used in page headers.
\title[Aggregation \& Analysis of Video Game Traces]{Research Proposal Asgn: Aggregation and Analysis of Video Game Traces for the Purposes of Archaeogaming}

%%
%% The "author" command and its associated commands are used to define
%% the authors and their affiliations.
%% Of note is the shared affiliation of the first two authors, and the
%% "authornote" and "authornotemark" commands
%% used to denote shared contribution to the research.
\author{Aly Cerruti}
\email{aly.cerruti@ucalgary.ca}
\affiliation{%
  \institution{University of Calgary}
  \city{Calgary}
  \state{Alberta}
  \country{Canada}
}

%%
%% By default, the full list of authors will be used in the page
%% headers. Often, this list is too long, and will overlap
%% other information printed in the page headers. This command allows
%% the author to define a more concise list
%% of authors' names for this purpose.
% \renewcommand{\shortauthors}{Trovato et al.}

%%
%% The abstract is a short summary of the work to be presented in the
%% article.
%\begin{abstract}
%  A clear and well-documented \LaTeX\ document is presented as an
%  article formatted for publication by ACM in a conference proceedings
%  or journal publication. Based on the ``acmart'' document class, this
%  article presents and explains many of the common variations, as well
%  as many of the formatting elements an author may use in the
%  preparation of the documentation of their work.
%\end{abstract}

%%
%% The code below is generated by the tool at http://dl.acm.org/ccs.cfm.
%% Please copy and paste the code instead of the example below.
%%
\begin{CCSXML}
    <ccs2012>
    <concept>
    <concept_id>10010405.10010476.10011187.10011190</concept_id>
    <concept_desc>Applied computing~Computer games</concept_desc>
    <concept_significance>300</concept_significance>
    </concept>
    <concept>
    <concept_id>10011007.10011074.10011111.10003465</concept_id>
    <concept_desc>Software and its engineering~Software reverse engineering</concept_desc>
    <concept_significance>500</concept_significance>
    </concept>
    <concept>
    <concept_id>10010405.10010432.10010434</concept_id>
    <concept_desc>Applied computing~Archaeology</concept_desc>
    <concept_significance>300</concept_significance>
    </concept>
    <concept>
    <concept_id>10003456.10003457.10003521.10003524</concept_id>
    <concept_desc>Social and professional topics~History of software</concept_desc>
    <concept_significance>300</concept_significance>
    </concept>
    </ccs2012>
\end{CCSXML}

\ccsdesc[300]{Applied computing~Computer games}
\ccsdesc[500]{Software and its engineering~Software reverse engineering}
\ccsdesc[300]{Applied computing~Archaeology}
\ccsdesc[300]{Social and professional topics~History of software}

%%
%% Keywords. The author(s) should pick words that accurately describe
%% the work being presented. Separate the keywords with commas.
\keywords{Archaeogaming, Program Traces, Time-travel Debugging}

% \received{20 February 2007}
% \received[revised]{12 March 2009}
% \received[accepted]{5 June 2009}

%%
%% This command processes the author and affiliation and title
%% information and builds the first part of the formatted document.
\maketitle

\section{Introduction}

\outline{
\begin{itemize}
  \item briefly introduce archaeogaming as a concept
  \item introduce the state of archaeogaming as a field
  \item introduce the tools needed for archaeogaming
\end{itemize}
}

\textit{Archaeogaming} is a relatively recent field focusing on the study of games under the lens of archaeology~\cite{reinhard2018archaeogaming}.
Many recent works in the field of archaeogaming specifically focus on video games and the machines that run them, specifically from a period of time where hand-written assembly was common.
For example, \cite{aycock2024experimental} by Aycock and Biittner demonstrates applied archaeogaming in the context of developing software using the Apple II for the Commodore 64.

This kind of archaeogaming warrants a few specialized tools.
Emulators are already showcased in~\cite{aycock2024experimental}, and MAME serves as a great starting point for much of archaeogaming, as demonstrated.
MAME includes a debugger, though the interface is very simple and doesn't support time-travel debugging (debugging where you can also ``step back'' instead of just forward).
However, of particular interest are traces: recordings of sessions of real human interaction that include every state and cycle of the underlying machine.
Traces allow one to do things such as compare the behavior of the same program between two different sessions, or compare the behavior of different programs that may have similar underlying systems.

\section{Motivation}

\outline{
\begin{itemize}
  \item motivate the need for tooling
  \item describe the current state of tooling
  \item motivate the new tool / research topic
\end{itemize}
}

Archaeogaming being a new field means it is lacking in the maturity in resources and maturity in tooling that other fields have.
One of archaeogaming's facets, reverse engineering, has plenty of tools like GDB, IDA, and Ghidra, that have been in development for decades, but aren't suited towards archaeogaming due to their focus on relatively modern environments.
This makes analyzing single games, not to mention a whole corpus of games together as one artifact, a relative challenge: these older machines can have memory mappers in external cartridges or features like banks and echo memory that completely sidestep traditional assumptions (namely, one virtual address mapping to one physical address).

These challenges lead to a need for new tools.
Specifically, a tool that allows the aggregation and analysis of a corpus of traces would answer many of these challenges.
One could iterate through individual states looking for certain conditions or find patterns among the whole dataset, among other analyses that may be useful for archaeogaming.
Some example usages for this tool (not exhaustive) may include:

\begin{description}
    \item[Single-Game Reverse Engineering] in which a user combines the SQL-like interface with a debugger in order to rewind and replay a specific section of code for the purpose of understanding that code segment.
    In modern settings, code is usually generated by a compiler, which makes the development of decompilers (such as those featured in IDA and Ghidra) an obvious solution.
    However, in our situation, the underlying binaries are effectively hand-written (having been assembled rather than compiled), and the architectures themselves make a representation in C much more difficult.
    \item[System Analysis] in which a user analyzes many traces over different games for the same system, looking for commonalities.
    For example, a user might want to analyze many different NES games and look for common behaviors.
    They might uncover a pattern of execution across games of ``wait for the zero-sprite interrupt, then change the scroll register'', which was a common method at the time of separating a HUD from a side-scrolling tile layer.
\end{description}

\section{Related Work}

\section{Methodology}

\outline{
\begin{itemize}
    \item start with a ``stupid'' analyzer in Python
    \item move to a ``less stupid'' SQL interface that uses a forward and backward copy of deltas
    \item move to a ``smart'' interface based on the ideas from PANDA
\end{itemize}
}

We may initially build such a tool as a prototype (i.e.\ in Python, with unoptimized logic) that scans traces generated by \textsc{Lava}.
\textsc{Lava} is an as-of-yet unpublished work, but an example of the trace format is shown in~\ref{fig:lava-snippet}.
As traces are incredibly large uncompressed, and scanning through them as files must be done linearly, this tool will be incredibly slow.
We will treat it as a baseline: the minimum required to be ``a tool that analyzes traces''.

\begin{figure}[ht]
    \begin{verbbox}
R	00cdfe4e	f45f	d5
R	00cdfe4e	f460	0d
R	00cdfe4e	ffff	00
R	00cdfe4e	d00d	40
D	00cdfe52	e040	f0f1	fee8	f14b	cbe8	d0	50	f461	BEQ    $F459
R	00cdfe52	f461	27
R	00cdfe52	f462	f6
R	00cdfe52	ffff	00
D	00cdfe55	e040	f0f1	fee8	f14b	cbe8	d0	50	f463	LDA    $C823
R	00cdfe55	f463	b6
R	00cdfe55	f464	c8
R	00cdfe55	f465	23
R	00cdfe55	ffff	00
R	00cdfe55	c823	01
D	00cdfe5a	0140	f0f1	fee8	f14b	cbe8	d0	50	f466	CLR    $0A
R	00cdfe5a	f466	0f
R	00cdfe5a	f467	0a
R	00cdfe5a	ffff	00
R	00cdfe5a	d00a	e0
R	00cdfe5a	f468	4d
W	00cdfe5a	d00a	00
I	SNAPSHOT 9 cdfe5a
D	00cdfe60	0140	f0f1	fee8	f14b	cbe8	d0	54	f468	TSTA
    \end{verbbox}
    \centering\theverbbox
    \caption{A snippet of a \textsc{Lava} trace. The \texttt{R} commands show memory reads, \texttt{W} writes, \texttt{D} disassembly, and \texttt{I} logging from the tracer.}
    \label{fig:lava-snippet}
\end{figure}

Considering we want an SQL-like interface as the frontend for developer familiarity, we may consider moving the approach towards actual usage of a SQL database.
We could compare the storage of traces singularly versus the storage of a forward trace and a reversed trace for faster or more efficient rewinding.

There are existing works regarding the storage of program traces in databases (\cite{andjelkovic2011trace},~\cite{montplaisir2013efficient},~\cite{dolan2015repeatable}) that will be important references for efficient storage of our traces.
It is likely incorporating ideas from these works will lead to further performance gains and insights on possible queries.

Regarding queries, we must carefully consider the interface provided:~\cite{osborn2015playspecs} and~\cite{goldsmith2005relational} detail the kinds of queries a programmer may want to make over state machines and program traces, respectively.
This will be done in tandem with incremental evaluation of the tool.

\section{Evaluation and Expectations}

\outline{
\begin{itemize}
    \item analyzing single games
    \item demonstrating reverse engineering usage
    \item analyzing a corpus of games (new)
    \item expecting to match the results of other reverse engineering works
    \item a new tool that enables analysis of systems and common practices
\end{itemize}
}

We expect to be able to match the results of other single-shot reverse engineering works, without a running copy of the program.
That is, we should be able to perform at least the same kinds of analysis using N traces with our tool as a regular reverse engineer would be able to perform with N live runs (without pokes).
This would be especially helpful to validate the use of traces themselves over live debugging for reverse engineering in the archaeogaming context.
We note that usage of the MAME debugger to record traces themselves is particularly slow, and direct time-travel is impossible by default in MAME\@.

We also expect the tool to prove its usefulness in the aforementioned areas of analysis.
We will demonstrate the single-game case this by including demonstrations of the tool being used for two or three significantly different games.
The utility will be mostly evident in the multi-game case: running a large set of programs at once under MAME, not to mention managing all of their debuggers at once, ends up being a headache.
Ideally our tool will be well-designed for the multi-game case to make analysis simple and practical.

We hope as a result of this work to fill one of the major gaps in tooling mentioned in the introduction.
It's worth remembering here that archaeogaming is an interdisciplinary field: not everyone interested in archaeogaming has a strong computer science background.
For this, the final major goal for the tool is to remain an open-source project with standards that prevent it from being classified under the ``academic code'' stereotype.
This is a goal that extends beyond the scope of the paper, and beyond the timeline of a 2-year MSc.

%%
%% The next two lines define the bibliography style to be used, and
%% the bibliography file.
\bibliographystyle{ACM-Reference-Format}
\ifcurrfileabsdir{}{
\bibliography{main}
}{
\bibliography{\currfileabsdir/main}
}

\end{document}
\endinput
%%
%% End of file `sample-manuscript.tex'.
